\chapter*{前言}
\addcontentsline{toc}{chapter}{前言}
“全等的两个三角形一定是相似的”，这一命题是正确的，那么我们就要加以严格证明；“相似的两个三角形一定是全等的”这一命题是不正确的，那么我们就要找出两个相似三角形并不全等，也就是说，举一个反例。由此看来，对于命题来说，给出证明和构造反例是同样重要的。

本书的主要内容是举出关于代数中的反例，包括高等代数和近世代数基础两部分。在前一部分中，环上的线性空间一节已超出高等代数的范围，但它对理解线性空间的结构不是无益的。在举反例的过程中，所涉及到的定理、命题的论证均可在高等代数和近世代数基础中找到。为了有助于对问题的理解，我们还加入了一些说明性的例题。

本书参考了张禾瑞教授著的《近世代数基础》及由王世强教授执笔的《关于代数教材的几点注记》的讲义等材料。

在编写过程中，张德荣同志精心审阅了全部手稿，提出许多改进意见，同时还得到我的老师和其他同志的帮助，谨此一并感谢。

希望同志们多加指正。

\begin{flushright}
1982 年 6 月
\end{flushright}
